\chapter{Introduzione}
\label{chap:introduzione}

\section{Contesto}

Un backup, nell'ambito della sicurezza informatica, è un processo di disaster recovery: un procedimento
che consente di recuperare dati andati persi, per esempio in seguito a un attacco, un guasto hardware o
software, o un errore umano. Una delle prime risposte tecniche alla perdita di dati è la ridondanza dei dischi su
cui questi sono memorizzati: nel 1988 Patterson, Gibson e Katz formalizzano il RAID (Redundant Array of
Inexpensive Disks), che distribuisce lo stesso dato su più dischi economici per ottenere un'affidabilità
paragonabile a quella di un singolo disco più costoso \cite{patterson-gibson-katz-raid-1988}.

Questa ridondanza, però, resta locale e non protegge da un evento che distrugga l'intero apparato. Da
qui nasce l'esigenza di un backup remoto: una copia dei dati conservata in un luogo diverso da quello
degli originali. Il modo più semplice è affidarli a terzi, un conoscente che possiede l'hardware o
un'attività commerciale che offre questo tipo di servizio, ma questo richiede di fidarsi che quei dati
non vengano letti, modificati o venduti da chi li conserva.

Chi offre il servizio, dal canto suo, trova conveniente servire più utenti contemporaneamente. Questo
richiede però di distinguere gli utenti, e di garantire che ciascuno acceda solo ai propri dati: in
pratica, richiede un sistema di autenticazione e controllo degli accessi.

Tra i diversi sistemi di autenticazione esistenti, quello adottato e discusso in questa tesi è basato su 
email e password che, in quanto dati sensibili, richiedono un livello di protezione elevata. 
Il principio su cui si fonda questo tipo di autenticazione è semplice:
l'utente si registra inviando al server le proprie credenziali; durante il login le invia di nuovo, e il
server le confronta con quelle salvate in fase di registrazione. Se coincidono, l'utente è autenticato.

Il fatto che il server possieda le credenziali espone quindi tutti gli utenti al rischio che vengano
rubate. Esistono numerosi sistemi per impedirlo; uno tra tutti, approfondito in questa tesi, è l'uso di
una funzione di hash.
Il server, al posto di salvare le credenziali in chiaro, le salva
sotto forma di hash. Durante il login, ricalcola gli hash delle credenziali fornite dall'utente e
verifica se coincidono con quelli salvati.

\subsection{Funzioni di hash}

Le funzioni di hash sono funzioni non invertibili che trasformano un input di lunghezza
arbitraria in un output di lunghezza fissa e strettamente correlato con i dati di input, chiamato hash.
Indicata con $H$ una funzione di hash, questa gode delle seguenti proprietà fondamentali:

\begin{itemize}
    \item \textbf{resistenza alla preimmagine:} dato $h$, è computazionalmente intrattabile trovare $x$
    tale che $H(x) = h$;
    \item \textbf{resistenza alla seconda preimmagine:} data $x$, è computazionalmente intrattabile
    trovare $y \neq x$ tale che $H(x) = H(y)$;
    \item \textbf{resistenza alle collisioni:} è computazionalmente intrattabile trovare una coppia
    $(x, y)$ con $x \neq y$ tale che $H(x) = H(y)$.
\end{itemize}


\section{Problema affrontato}

Tra le cause di perdita dei dati elencate sopra, RAM-USB si concentra sull'attacco da parte di un 
malintenzionato, affidandosi ai principi zero-trust e defense-in-depth nel tentativo di garantire la 
sicurezza dei dati anche quando un componente qualsiasi del sistema viene compromesso.

\section{Obiettivi}

RAM-USB, Remotely Accessible Multi-User Secure Backup è un caso di studio approfondito sulla progettazione, 
implementazione e distribuzione di un servizio di backup che soddisfi i requisiti
descritti nel \autoref{chap:requisiti}. Non è pensato per competere con le soluzioni 
commerciali esistenti, ma per discuterne le scelte di design, illustrandone pro e contro in modo 
trasparente, con un approccio security first e concentrato sulla difesa del sistema da attacchi esterni.


\section{Metodologia}
Parte dei componenti di RAM-USB nascono dalla fase di brainstorming condotta insieme a Riccardo 
Gottardi \cite{gottardi-github}.

La metodologia adottata combina il modello a spirale per la definizione e l'evoluzione dei 
requisiti \cite{Boehm88} con l'approccio Test-Driven Development per l'implementazione e il modello a V per 
la loro verifica \cite{ramusb-contributing}.

Il linguaggio usato per l'intero stack è Go, un linguaggio compilato progettato esplicitamente per gestire 
in modo efficiente molte connessioni di rete simultanee.

Per orchestrare l'infrastruttura attorno al codice Go (avvio dei servizi, supervisione dei processi nei 
container, sequenziamento all'avvio) è stato usato il linguaggio Bash. 

Il codice è versionato con Git e pubblicato gratuitamente, in versione integrale su GitHub, distribuito con 
licenza MIT. 

L'implementazione è stata interamente realizzata con l'assistenza dei Large Language Model Claude 
Sonnet 4.5 e Sonnet 5, in forma multi-agentica tramite l'harness Claude Code eseguito in una macchina 
virtuale su Proxmox per questioni di sicurezza e privacy.

L'intelligenza artificiale generativa è stata impiegata esclusivamente come strumento sotto stretta 
supervisione umana. 
La progettazione, la struttura del codice, i confini dei moduli, la firma, il comportamento delle singole 
funzioni, e i test corrispondenti restano una decisione umana; l'assistente ha implementato solo quanto 
specificato, in passi piccoli, incrementali e verificabili, prendendosi la libertà implementativa di basso 
livello.
Ogni modifica è stata revisionata ed esplicitamente approvata prima di essere introdotta nella codebase.
Le numerose strategie di mitigazione degli errori implementativi e allucinazioni dei due modelli, nonchè la 
getione dei multipli agenti utilizzati non sono oggetto di questa tesi. 

\section{Struttura della tesi}

Il \autoref{chap:principi-fondanti} approfondisce i due principi di sicurezza su cui RAM-USB è costruito,
zero-trust e defense-in-depth, dalle loro origini alla loro applicazione all'informatica.

Il \autoref{chap:requisiti} descrive la metodologia con cui i requisiti sono stati raccolti e tracciati,
e i casi d'uso che ne derivano.

Il \autoref{chap:architettura} presenta i componenti che partecipano alla registrazione e al login.

Il \autoref{chap:autenticazione} segue le credenziali lungo quel percorso, dal confine pubblico fino al
record salvato, e descrive le regole che ogni strato applica.

Il \autoref{chap:storage} segue il backup fino al proprio spazio isolato su Storage-Service.

Il \autoref{chap:rete} descrive come la rete mesh applichi i confini di raggiungibilità tra i
componenti.

Il \autoref{chap:osservabilita} chiude il ciclo di vita di una richiesta con la pipeline di
osservabilità che la rende misurabile.

Il \autoref{chap:testing} verifica il sistema così costruito attraverso i quattro livelli del modello a
V, dal test di unità alla lista di accettazione, e riporta cosa ha trovato eseguendolo su hardware reale.

Il \autoref{chap:deployment} descrive come questi componenti siano effettivamente distribuiti, tra
container Docker, macchine Proxmox e un VPS pubblico.

Il \autoref{chap:sicurezza} rilegge l'intero sistema dal punto di vista di chi tenta di comprometterlo,
verificando quali garanzie resistano anche alla sua compromissione totale.

Il \autoref{chap:conclusioni} chiude la tesi bilanciando i risultati ottenuti con i limiti riscontrati,
e indicando gli sviluppi futuri più rilevanti.
